\pdfminorversion=7
\documentclass[11pt,letterpaper]{article}

\usepackage[utf8]{inputenc}
\usepackage[T1]{fontenc}
\usepackage[french]{babel}
\frenchsetup{SmallCapsFigTabCaptions=false}
\usepackage[scale=0.95]{tgheros}
\renewcommand{\familydefault}{\sfdefault}
\usepackage{geometry}
\usepackage{graphicx}
\usepackage{grffile}
\usepackage{float}
\usepackage{booktabs}
\usepackage{longtable}
\usepackage{array}
\usepackage{tabularx}
\usepackage[table]{xcolor}
\usepackage{tikz}
\usetikzlibrary{babel,positioning,arrows.meta,fit,calc,shapes.geometric}
\usepackage{eso-pic}
\usepackage{fancyhdr}
\usepackage{titlesec}
\usepackage[font=small,labelfont=bf]{caption}
\usepackage{microtype}
\usepackage{indentfirst}
\usepackage{enumitem}
\usepackage[most]{tcolorbox}
\usepackage{hyperref}

\graphicspath{{assets/}{img/}}

\geometry{
    letterpaper,
    left=0.95in,
    right=0.95in,
    top=1.55in,
    bottom=0.78in,
    headheight=18pt,
    headsep=18pt,
    footskip=28pt
}

\definecolor{tameoBlue}{HTML}{173B5C}
\definecolor{tameoLink}{HTML}{3A6B78}
\definecolor{tameoText}{HTML}{13283A}
\definecolor{tablerow}{RGB}{242, 244, 244}
\definecolor{tameoPale}{HTML}{E8EEF3}
\definecolor{okGreen}{HTML}{2E7D4F}
\definecolor{warnAmber}{HTML}{B26A00}
\definecolor{openRed}{HTML}{A33A3A}
\newcommand{\tameoCoverIndent}{0.20in}

% ---------------------------------------------------------------------------
% Informations du document
% ---------------------------------------------------------------------------
\newcommand{\tameoTitle}{Bateau IA -- Groupe 25}
\newcommand{\tameoDocumentType}{Résumé exécutif -- Système de détection et d'interception d'obstacles marins}
\newcommand{\tameoPdfTitle}{Résumé exécutif TAMEO -- Détection et interception d'obstacles marins}
\newcommand{\tameoPdfAuthor}{Projet TAMEO}

\newcommand{\tameoAuthorsBlockA}{%
    \tameoAuthorCell{CAVALLARO}{Vinicius Cesar}{vinicius-cesar.cavallaro@ensta.fr}
    & \tameoAuthorCell{COUERRE}{Loup}{loup.couerre@ensta.fr}
    \tabularnewline[0.56in]
    \tameoAuthorCell{DIAS LEITE}{Julia Ellen}{julia-ellen.dias@ensta.fr}
    & \tameoAuthorCell{MARCHÉ}{Benoît}{benoit.marche@ensta.fr}
}

\newcommand{\tameoAuthorsBlockB}{%
    \tameoAuthorCell{DA SILVEIRA}{Lucas}{lucas.da-silveira-absalao@ensta.fr}
    & \tameoAuthorCell{DE SOUZA AMODIO}{Lucca}{lucca.de-souza-amodio@ensta.fr}
    \tabularnewline[0.56in]
    \tameoAuthorCell{NOUKOUA TATMFO}{Maeva sandy}{maeva-sandy.noukoua@ensta.fr}
    & \tameoAuthorCell{TARAZONA JIMENEZ}{Javier andres}{javier-andres.tarazona@ensta.fr}
}

\hypersetup{
    colorlinks=true,
    linkcolor=tameoBlue,
    urlcolor=tameoLink,
    pdftitle={\tameoPdfTitle},
    pdfauthor={\tameoPdfAuthor}
}

\setlength{\parindent}{1.5em}
\setlength{\parskip}{0.35em}
\renewcommand{\arraystretch}{1.18}
\newcolumntype{L}[1]{>{\raggedright\arraybackslash}p{#1}}
\newcolumntype{Y}{>{\raggedright\arraybackslash}X}
\setlength{\emergencystretch}{1.5em}
\setlist{itemsep=0.15em,topsep=0.3em}

\titleformat{\section}
  {\sffamily\bfseries\LARGE\color{tameoBlue}}
  {\thesection}{0.75em}{}
\titleformat{\subsection}
  {\sffamily\bfseries\large\color{tameoBlue}}
  {\thesubsection}{0.75em}{}
\titlespacing*{\section}{0pt}{1.3em}{0.65em}
\titlespacing*{\subsection}{0pt}{1.05em}{0.35em}

\pagestyle{fancy}
\fancyhf{}
\renewcommand{\headrulewidth}{0pt}
\fancyfoot[R]{\thepage}

\newcommand{\tameoBranding}{%
    \begin{tikzpicture}[remember picture,overlay]
        \node[anchor=north east] at ([xshift=-1.16in,yshift=-0.58in]current page.north east)
            {\includegraphics[width=1.98in]{tameo-logo.png}};
        \node[anchor=north east] at ([xshift=-1.18in,yshift=-1.12in]current page.north east)
            {\includegraphics[width=2.05in,height=0.026in]{tameo-bar.png}};
    \end{tikzpicture}%
}

\newcommand{\tameoAuthorCell}[3]{%
    \begin{minipage}[t]{2.6in}
        \raggedright
        \hyphenpenalty=10000
        \exhyphenpenalty=10000
        {\sffamily\fontsize{10.8}{12.6}\selectfont #1\par}
        {\sffamily\fontsize{10.2}{12}\selectfont #2\par}
        {\sffamily\fontsize{6.45}{7.6}\selectfont\href{mailto:#3}{#3}\par}
    \end{minipage}%
}

\newcommand{\maketameotitle}{%
    \begin{titlepage}
        \thispagestyle{empty}
        \begin{tikzpicture}[remember picture,overlay]
            \node[anchor=north east] at ([xshift=-1.37in,yshift=-2.42in]current page.north east)
                {\includegraphics[width=1.82in]{tameo-logo.png}};
            \node[anchor=north west] at ([xshift=1.16in,yshift=-3.00in]current page.north west)
                {\includegraphics[width=6.55in,height=0.115in]{tameo-bar.png}};
        \end{tikzpicture}

        \vspace*{1.82in}
        {\noindent\hspace*{\tameoCoverIndent}\sffamily\bfseries\fontsize{42}{50}\selectfont\color{tameoBlue}\tameoTitle\par}
        \vspace{0.45in}
        {\noindent\hspace*{\tameoCoverIndent}\parbox{5.9in}{\raggedright\sffamily\itshape\fontsize{16}{20}\selectfont\color{tameoText}\tameoDocumentType}\par}

        \vspace{0.46in}
        \noindent\hspace*{\tameoCoverIndent}\makebox[0pt][l]{%
            \begin{tabular*}{5.95in}{@{\extracolsep{\fill}}ll@{}}
                \tameoAuthorsBlockA
            \end{tabular*}%
        }\par
        \vspace{0.25in}
        \noindent\hspace*{\tameoCoverIndent}\makebox[0pt][l]{%
            \begin{tabular*}{5.95in}{@{\extracolsep{\fill}}ll@{}}
                \tameoAuthorsBlockB
            \end{tabular*}%
        }
    \end{titlepage}
}

% Encadrés et indicateurs propres au résumé exécutif
\newtcolorbox{tameoBox}[1][]{
    enhanced, breakable,
    colback=tameoPale, colframe=tameoBlue,
    boxrule=0pt, leftrule=3pt, arc=0pt,
    left=8pt, right=8pt, top=6pt, bottom=6pt,
    fonttitle=\sffamily\bfseries\color{tameoBlue},
    coltitle=tameoBlue, colbacktitle=tameoPale,
    title={#1}, attach title to upper={\par\smallskip}
}

\newcommand{\kpi}[2]{%
    \begin{minipage}[t]{0.235\linewidth}
        \centering
        {\sffamily\bfseries\fontsize{19}{22}\selectfont\color{tameoBlue}#1\par}
        \vspace{2pt}
        {\sffamily\fontsize{8.4}{10}\selectfont\color{tameoText}#2\par}
    \end{minipage}%
}

\newcommand{\statut}[1]{%
    \ifx#1v{\color{okGreen}\textbf{Validé}}\fi%
    \ifx#1p{\color{warnAmber}\textbf{Partiel}}\fi%
    \ifx#1o{\color{openRed}\textbf{Ouvert}}\fi%
}

\begin{document}

\maketameotitle
\setcounter{page}{2}
\AddToShipoutPictureBG{\tameoBranding}

% ===========================================================================
\section{Synthèse}
% ===========================================================================

\begin{tameoBox}[En une phrase]
Pour sa première participation à la classe IA du Monaco Energy Boat Challenge
(MEBC, 8--11 juillet 2026), l'association TAMEO et l'ENSTA Paris ont développé
un système embarqué qui \textbf{détecte les objets marins} (bouées fixes, bouée
mobile, quai), \textbf{les localise} dans le plan d'eau et \textbf{décide,
sans intervention humaine, de les éviter ou de les intercepter} (atteindre,
contourner, accoster), sous la supervision d'un pilote embarqué.
\end{tameoBox}

\medskip
\noindent
\kpi{0,993}{mAP50 du meilleur modèle YOLO26 (deux classes de bouées)}\hfill
\kpi{0,888}{mAP50-95, critère plus strict sur le même jeu de test}\hfill
\kpi{50 m}{portée maximale de détection configurée sur les caméras ZED}\hfill
\kpi{5 s}{horizon de prédiction du planificateur local DWA}
\medskip

\noindent\textbf{Ce qui a été réalisé.}
\begin{itemize}
    \item Une \textbf{chaîne de perception complète} : constitution d'un jeu de
    données maritime, annotation CVAT, entraînement YOLO26, export ONNX et
    inférence sur Jetson AGX Orin avec caméras stéréo ZED X One.
    \item Une \textbf{boucle de décision} fondée sur la \emph{Dynamic Window
    Approach} (DWA) : à chaque pas, le bateau évalue les commandes de vitesse
    et de lacet physiquement atteignables, prédit leurs trajectoires et retient
    celle qui progresse vers l'objectif en restant loin des obstacles et dans
    la zone autorisée.
    \item Des \textbf{primitives d'interception} pour les épreuves où la cible
    n'a pas de position connue : sélection de la bouée la plus proche, balayage
    de recherche, critère de contournement, accostage en deux temps
    (alignement puis entrée).
    \item Une \textbf{chaîne de commande sécurisée} vers les actionneurs, avec
    arbitrage manuel/IA, signalisation par LED et repli en cas de perte de la
    radiocommande.
\end{itemize}

\noindent\textbf{Ce qui reste à démontrer.} Les performances de détection sont
mesurées sur le jeu de données, pas encore sur l'eau à Monaco. La portée réelle,
la latence de bout en bout, le comportement face à la bouée mobile (aujourd'hui
traitée comme un obstacle fixe, mis à jour à chaque image) et le réglage
dynamique du bateau réel restent les principaux points ouverts.

\noindent\textbf{Recommandation.} Concentrer la prochaine phase sur une
campagne de mesures sur l'eau (portée, latence, faux positifs), la prédiction
de trajectoire de la bouée mobile et le réglage des limites dynamiques du
bateau réel, avant toute recherche de performance pure.

% ===========================================================================
\section{Contexte et enjeux}
% ===========================================================================

Le MEBC, organisé par le Yacht Club de Monaco, ouvre une \textbf{classe IA}
pour des bateaux autonomes de niveau 5 (définition SAE) : pendant une course,
le bateau doit être piloté uniquement par des algorithmes, sans télécommande
(toute manipulation à distance entraîne la disqualification). Un pilote reste
à bord pour sélectionner l'épreuve, superviser et reprendre la main en cas
d'urgence. Quatre épreuves sont au programme : \emph{Figure 8},
\emph{Slalom}, \emph{Docking} et \emph{Sprint}.

Le projet a été mené par l'équipe PIE \og Bateau IA \fg{} de l'ENSTA Paris
avec les étudiants de première année de TAMEO. Deux contraintes ont façonné
la solution : l'\textbf{accès tardif au bateau final} (un Hobie Cat 14 sur
châssis aluminium, au gabarit mesuré d'environ 5\,m $\times$ 2,9\,m), compensé par l'usage d'un
BlueBoat de Blue Robotics comme banc d'essai, et la \textbf{rareté des images
de bouées} exploitables, compensée par une démarche centrée sur les données.
Le budget global de l'association est maintenu à 50\,481\,\texteuro{}.

L'enjeu du système décrit ici est double : \textbf{sécurité} (ne heurter ni
bouée, ni quai, ni yacht amarré) et \textbf{performance} (contourner au plus
vite les bouées imposées). Les deux reposent sur la même capacité : savoir ce
qui se trouve devant le bateau, où, et quoi en faire.

% ===========================================================================
\section{Le système de détection et d'interception}
% ===========================================================================

\subsection{Architecture fonctionnelle}

Le système suit une boucle \emph{percevoir -- localiser -- décider --
commander -- agir}, encadrée en permanence par une couche de supervision et de
sécurité (figure~\ref{fig:chaine}).

\begin{figure}[H]
\centering
\begin{tikzpicture}[
    font=\sffamily,
    hdr/.style={fill=tameoBlue, text=white, font=\sffamily\bfseries\scriptsize,
                minimum width=2.75cm, minimum height=0.55cm, align=center, inner sep=2pt},
    body/.style={draw=tameoBlue, line width=0.8pt, fill=white, text width=2.55cm,
                 minimum width=2.75cm, minimum height=2.0cm, align=center, inner sep=3pt,
                 font=\sffamily\scriptsize},
    arr/.style={-{Stealth[length=2.4mm]}, line width=0.9pt, color=tameoBlue},
    sarr/.style={-{Stealth[length=2mm]}, line width=0.6pt, color=tameoLink, dashed}
]
\foreach \i/\t/\b in {
    1/PERCEVOIR/{Caméras ZED X One\\ YOLO26 (ONNX)\\ suivi des objets\\ portée 50\,m},
    2/LOCALISER/{Profondeur stéréo\\ repères caméra,\\ bateau, monde\\ GPS/RTK, IMU},
    3/DÉCIDER/{Planificateur DWA\\ \textbf{éviter} : coût obstacle\\ \textbf{intercepter} : cible, contournement},
    4/COMMANDER/{Consignes $v$, $\omega$\\ MAVLink vers le CPN\\ arbitrage par le GAC},
    5/AGIR/{Propulsion (VESC)\\ safran (ODrive)\\ actionneurs à 0\textdegree{} en fin de course}
}{
    \node[hdr] (h\i) at ({(\i-1)*3.2},0) {\t};
    \node[body, anchor=north] (b\i) at (h\i.south) {\b};
}
\foreach \i/\j in {1/2,2/3,3/4,4/5}{ \draw[arr] (b\i.east) -- (b\j.west); }

\node[draw=tameoLink, fill=tameoPale, line width=0.6pt,
      text width=15.2cm, minimum height=0.9cm, align=center, font=\sffamily\scriptsize]
      (sup) at ($(b3.south)+(0,-1.2)$)
      {\textbf{SUPERVISER ET SÉCURISER}\\[1pt] pilote à bord (AMS, kill cord, reprise en main)
       \,$\cdot$\, LED d'état \,$\cdot$\, dashboard IA \,$\cdot$\, repli si perte radio $>$ 500\,ms};
\foreach \i in {1,3,4,5}{ \draw[sarr] (sup.north -| b\i.south) -- (b\i.south); }

\node[font=\sffamily\scriptsize\itshape, text=tameoText, anchor=south west] (env)
      at ($(h1.north west)+(0,0.55)$) {Scène marine : bouées, quai, yachts, reflets};
\draw[arr] (env.south -| h1.north) -- (h1.north);
\draw[sarr] (h5.north) -- ++(0,0.4) -| node[pos=0.25, above, font=\sffamily\scriptsize\itshape, text=tameoText]
      {état du bateau (position, cap, vitesse)} (h2.north);
\end{tikzpicture}
\caption{Chaîne fonctionnelle du système de détection et d'interception.}
\label{fig:chaine}
\end{figure}

\begin{itemize}
    \item \textbf{Percevoir.} Les caméras ZED X One (liaison GMSL2) alimentent
    une Jetson AGX Orin 64\,Go. Un modèle YOLO26 affiné sur deux classes,
    \texttt{buoy\_dynamic} et \texttt{buoy\_fixed}, est exécuté par le SDK
    ZED, qui ajoute la profondeur stéréo et le suivi des objets (seuil de
    confiance de 50\,\%). Pour le docking, un modèle de segmentation distingue
    l'eau du quai ; un second modèle, plus léger, n'est sollicité que si le
    premier manque de confiance, au plus une fois par seconde, pour préserver
    la latence.
    \item \textbf{Localiser.} Chaque détection est convertie du repère caméra
    au repère bateau puis aux coordonnées géographiques à l'aide du GPS (RTK
    issu d'un service licencié, comme l'impose le règlement). Elle est publiée
    avec sa classe, sa confiance, son identifiant de suivi et son rayon dans un
    état partagé, accessible à tous les modules.
    \item \textbf{Décider.} Le planificateur DWA échantillonne les commandes
    atteignables (pas de 0,1\,s, horizon de 5\,s) et les note selon cinq
    critères : progression vers l'objectif, vitesse, distance aux obstacles,
    proximité du but et douceur de la commande, auxquels s'ajoute le respect
    du polygone de la zone de course. Les obstacles sont la somme des bouées
    connues du parcours et des bouées détectées en direct.
    \item \textbf{Commander et agir.} Les consignes de vitesse et de vitesse de
    lacet partent en MAVLink vers le contrôleur principal de navigation (CPN :
    Raspberry Pi~4 et carte Navigator sous BlueOS/ArduPilot), qui dialogue par
    scripts Lua avec le gestionnaire automatique des commandes (GAC : Arduino
    Due). Le GAC arbitre entre commandes manuelles et IA, puis pilote le
    moteur de propulsion et le safran électrique.
\end{itemize}

\subsection{Éviter ou intercepter : une même boucle, deux intentions}

Le terme \og obstacle \fg{} recouvre deux situations opposées. Une bouée peut
être un \textbf{danger à éviter} (bouée mobile du slalom, yachts, bouées
étrangères au parcours) ou une \textbf{cible à intercepter} (marque de virage
du sprint ou du docking, dont la position est inconnue). Le système traite les
deux avec la même perception et le même planificateur ; seule change la
fonction de coût, ou l'objectif fixé au DWA. La logique du
tableau~\ref{tab:interception} est implémentée pour le docking ; celle du
sprint, conçue sur les mêmes principes, reste à coder.

\begin{table}[H]
\centering
\caption{Logique d'interception d'une cible de position inconnue}
\label{tab:interception}
\footnotesize
\rowcolors{2}{tablerow}{white}
\begin{tabularx}{\textwidth}{L{2.6cm}YL{4.3cm}}
\toprule
\textbf{Étape} & \textbf{Comportement} & \textbf{Paramètres actuels} \\
\midrule
1. Acquisition & Choix de la bouée détectée la plus proche ; priorité à la bouée mobile, puis aux bouées fixes. & Mise à jour à chaque image. \\
2. Recherche & Si rien n'est vu : rotation sur place de part et d'autre du cap initial. & $\pm$60\textdegree{} ; à défaut, retour au point de départ. \\
3. Approche & Le DWA prend la bouée comme but tout en évitant les autres objets. & Vitesse plafonnée à 2\,m/s. \\
4. Contournement & Validé quand le bateau a tourné d'au moins un demi-tour autour de la bouée en s'en approchant suffisamment, quel que soit le sens. & 180\textdegree{} cumulés, distance minimale $\leq$ 6\,m. \\
5. Accostage (docking) & Recherche du quai, alignement sur l'axe de la place, puis entrée lente en ligne droite et arrêt. & 0,4\,m/s ; arrêt à 3\,m de l'antenne GPS, soit environ 0,5\,m de l'étrave. \\
\bottomrule
\end{tabularx}
\end{table}

\subsection{Maturité du système}

\begingroup
\footnotesize
\rowcolors{2}{tablerow}{white}
\begin{longtable}{L{4.1cm}L{1.35cm}L{9.6cm}}
\caption{État de maturité des briques de détection et d'interception}\label{tab:maturite}\\
\toprule
\textbf{Brique} & \textbf{Statut} & \textbf{Constat} \\
\midrule
\endfirsthead
\toprule
\textbf{Brique} & \textbf{Statut} & \textbf{Constat} \\
\midrule
\endhead
Détection des bouées & \statut{v} & mAP50 de 0,993 et rappel de 0,970 sur le jeu de test ; 467 prédictions à 0,93 de confiance moyenne. Mesure hors de l'eau. \\
Intégration Jetson et caméras & \statut{v} & Deux flux ZED 1920\,$\times$\,1200 avec détection ONNX opérationnels en mode déploiement (juillet 2026). \\
Détections utilisées comme obstacles & \statut{v} & Les bouées vues en direct alimentent le DWA du bateau final, ce qui n'était pas encore le cas sur la version BlueBoat. \\
Segmentation du quai & \statut{p} & Modèles entraînés et intégrés ; aucune métrique validée en conditions réelles. \\
Bouée mobile du slalom & \statut{p} & Évitée comme un point fixe rafraîchi à chaque image : réactif, pas anticipatif. \\
Interception au docking & \statut{p} & Recherche, approche, contournement et accostage codés et jouables en simulation ; réglages provisoires, à recaler sur l'eau. \\
Interception au sprint & \statut{o} & Algorithme de recherche de la marque conçu, pas encore implémenté. \\
Recherche active d'une cible & \statut{o} & Au-delà du balayage sur place, fonction prévue mais non implémentée ; un échec déclenche le retour au départ. \\
Fusion des doublons & \statut{o} & Une même bouée vue deux fois peut produire deux obstacles. \\
Dynamique du bateau réel & \statut{o} & Limites de vitesse et de lacet héritées du BlueBoat, à recaler sur la coque de 5\,m. \\
Latence et portée réelles & \statut{o} & Non mesurées ; cibles du cahier des charges : vision $<$ 100\,ms, boucle $<$ 300\,ms, bouée vue à 20\,m. \\
\bottomrule
\end{longtable}
\endgroup

% ===========================================================================
\section{Périmètre, parties prenantes, environnement et missions}
% ===========================================================================

Cette section pose le cadre d'ingénierie système du \textbf{système
d'intérêt} : le système de détection et d'interception d'obstacles marins
(SdI) embarqué sur le bateau IA de TAMEO.

\subsection{Périmètre du système}

\begin{table}[H]
\centering
\caption{Frontière du système d'intérêt}
\label{tab:perimetre}
\footnotesize
\begin{tabularx}{\textwidth}{YYY}
\toprule
\textbf{Dans le périmètre} & \textbf{À la frontière (interfaces)} & \textbf{Hors périmètre} \\
\midrule
\textbullet~Caméras de perception et leur calibration \newline
\textbullet~Modèles de détection et de segmentation, inférence embarquée \newline
\textbullet~Localisation des objets (caméra, bateau, monde) \newline
\textbullet~Planification locale (DWA) et logique d'évitement/interception \newline
\textbullet~Logique de mission par épreuve \newline
\textbullet~Dashboard de visualisation IA \newline
\textbullet~Replis logiciels (perte de cible, retour au départ)
&
\textbullet~GPS/RTK et centrale inertielle \newline
\textbullet~Liaison MAVLink vers le CPN (consignes $v$, $\omega$ ; messages \texttt{GO\_RACE}, \texttt{RACE\_MODE}, \texttt{ERROR}) \newline
\textbullet~GAC : modes manuel/IA, LED, radiocommande \newline
\textbullet~Interrupteur AMS et sélection de l'épreuve par le pilote \newline
\textbullet~Télémétrie de sécurité vers la terre ($\geq$ 2 NM)
&
\textbullet~Coque, châssis et flottabilité \newline
\textbullet~Propulsion, batterie et distribution électrique \newline
\textbullet~Chaîne d'arrêt matérielle (kill cord, arrêt d'urgence) \newline
\textbullet~Équipements de sécurité du pilote \newline
\textbullet~Bateau de la classe Energy \newline
\textbullet~Logistique et transport vers Monaco \\
\bottomrule
\end{tabularx}
\end{table}

Le choix de frontière est délibéré : la chaîne d'arrêt d'urgence reste
\textbf{matérielle et indépendante} du SdI, de sorte qu'une défaillance
logicielle ne puisse jamais empêcher l'arrêt du bateau.

\subsection{Parties prenantes à intégrer}

\begin{figure}[H]
\centering
\begin{tikzpicture}[
    font=\sffamily\scriptsize,
    soi/.style={ellipse, draw=tameoBlue, line width=1.1pt, fill=tameoPale,
                minimum width=4.4cm, minimum height=1.9cm, align=center,
                font=\sffamily\bfseries\small, text=tameoBlue},
    act/.style={draw=tameoLink, line width=0.6pt, rounded corners=2pt, fill=white,
                text width=3.3cm, align=center, inner sep=4pt},
    lnk/.style={{Stealth[length=1.8mm]}-{Stealth[length=1.8mm]}, line width=0.6pt, color=tameoLink},
    one/.style={-{Stealth[length=1.8mm]}, line width=0.6pt, color=tameoLink},
    lbl/.style={font=\sffamily\scriptsize, text=tameoText, fill=white, inner sep=1.5pt, align=center}
]
\node[soi] (s) at (0,0) {Système de détection\\et d'interception};

\node[act] (env) at (-5.6, 3.1) {\textbf{Environnement marin}\\bouées, quai, yachts, houle, reflets};
\node[act] (rtk) at ( 0.0, 3.5) {\textbf{Service GNSS/RTK}\\licencié};
\node[act] (mebc) at ( 5.6, 3.1) {\textbf{Comité technique et\\ direction de course MEBC}};
\node[act] (sol) at ( 6.0, 0.0) {\textbf{Équipe à terre}\\station sol, télémétrie};
\node[act] (dash) at ( 5.6,-3.1) {\textbf{Jury, public}\\dashboard IA, diffusion};
\node[act] (plat) at ( 0.0,-3.5) {\textbf{Plateforme bateau}\\CPN, GAC, moteur, safran};
\node[act] (pil) at (-5.6,-3.1) {\textbf{Pilote embarqué}\\superviseur de sécurité};
\node[act] (team) at (-6.0, 0.0) {\textbf{TAMEO et ENSTA}\\conception, essais,\\ maintenance};

\draw[one] (env) -- node[lbl, sloped]{images, profondeur} (s);
\draw[one] (rtk) -- node[lbl]{position corrigée} (s);
\draw[lnk] (mebc) -- node[lbl, sloped]{règles, inspections, arrêt IA} (s);
\draw[lnk] (sol) -- node[lbl, above=1pt, fill=none]{sécurité} (s);
\draw[one] (s) -- node[lbl, sloped]{détections, trajectoire} (dash);
\draw[lnk] (plat) -- node[lbl]{consignes ($v$, $\omega$), état} (s);
\draw[lnk] (pil) -- node[lbl, sloped]{épreuve, AMS, reprise} (s);
\draw[lnk] (team) -- node[lbl, above=1pt, fill=none]{modèles} (s);
\end{tikzpicture}
\caption{Diagramme de contexte du système d'intérêt et de ses parties prenantes.}
\label{fig:contexte}
\end{figure}

\begingroup
\footnotesize
\rowcolors{2}{tablerow}{white}
\begin{longtable}{L{3.0cm}L{3.5cm}L{5.3cm}L{2.2cm}}
\caption{Parties prenantes et attentes vis-à-vis du système}\label{tab:pp}\\
\toprule
\textbf{Partie prenante} & \textbf{Rôle} & \textbf{Attentes et exigences clés} & \textbf{Intégration} \\
\midrule
\endfirsthead
\toprule
\textbf{Partie prenante} & \textbf{Rôle} & \textbf{Attentes et exigences clés} & \textbf{Intégration} \\
\midrule
\endhead
TAMEO & Client final, exploitant, porteur du projet & Bateau compétitif et sûr, réutilisable d'une année sur l'autre & Décideur \\
ENSTA Paris & Tutelle pédagogique, financeur & Démarche d'ingénierie rigoureuse, livrables documentés & Décideur \\
Yacht Club de Monaco, comité technique MEBC & Organisateur, autorité réglementaire & Conformité au règlement technique : aucune télécommande, AMS, LED, inspection à terre puis sur l'eau & Prescripteur \\
Direction de course & Arbitre pendant les épreuves & Pouvoir exiger l'arrêt de l'IA à tout moment ; autorise ou non une relance & Prescripteur \\
Pilote embarqué & Superviseur de sécurité & Choisir l'épreuve, comprendre l'état du système (LED, codes d'erreur), reprendre la main instantanément & Utilisateur \\
Équipe à terre & Supervision, communication avec le pilote & Données de sécurité en temps réel à plus de 2 NM, sans Wi-Fi & Utilisateur \\
Équipes PIE et étudiants 1A & Concepteurs, intégrateurs, testeurs & Interfaces claires entre vision, navigation et embarqué ; environnement reproductible & Contributeur \\
Sponsors (DGA, SaveMyBattery) & Financement, composants & Visibilité, usage pertinent du matériel fourni & Contributeur \\
Fournisseurs technologiques & Stereolabs, NVIDIA, Blue Robotics, ArduPilot, Ultralytics & Compatibilité des versions (JetPack, CUDA, SDK ZED) : dépendance critique & Contributeur \\
Usagers du plan d'eau & Yachts amarrés, autres concurrents, trafic résiduel & Ne subir aucun contact ni comportement imprévisible & Impacté \\
Jury et public & Évaluation, prix de l'innovation, diffusion & Dashboard lisible : boîtes de détection, distances, trajectoire planifiée & Impacté \\
Promotions futures & Repreneurs du système & Code, données et procédures transmissibles & Impacté \\
\bottomrule
\end{longtable}
\endgroup

\subsection{Environnement opérationnel}

\begingroup
\footnotesize
\rowcolors{2}{tablerow}{white}
\begin{longtable}{L{2.6cm}L{5.6cm}L{5.8cm}}
\caption{Environnement du système et conséquences de conception}\label{tab:env}\\
\toprule
\textbf{Dimension} & \textbf{Caractéristiques} & \textbf{Conséquences pour le système} \\
\midrule
\endfirsthead
\toprule
\textbf{Dimension} & \textbf{Caractéristiques} & \textbf{Conséquences pour le système} \\
\midrule
\endhead
Site & Port Hercule et stade nautique de Monaco (bouées vers 43\textdegree{}44,2'\,N -- 7\textdegree{}25,7'\,E) ; bassin bordé de yachts amarrés ; trafic réduit au minimum & Polygone de navigation autorisé dans le planificateur ; carte OpenStreetMap du port pour la planification \\
Conditions météo-marines & Mer ouverte jusqu'à force 6 Beaufort et 2\,m de vagues (catégorie C exigée pour le Sprint) ; embruns, sel, chaleur estivale & Électronique en boîtier IP56 refroidi ; tangage qui perturbe la caméra et la profondeur \\
Conditions visuelles & Soleil rasant ou zénithal, reflets sur l'eau, éblouissement, flou de mouvement & Augmentations de données ciblées ; risque de faux positifs et de bouées manquées \\
Objets présents & Bouées fixes, une bouée mobile lente (slalom), bouées de tailles et de formes variées autour du stade (sprint), quai et pontons & Deux classes de bouées distinctes ; discrimination de la cible parmi des bouées étrangères au parcours \\
Cadre réglementaire & Autonomie de niveau 5 ; aucune commande à distance ; vitesse $\geq$ 3 nœuds hors docking ; actionneurs à 0\textdegree{} avant et après la course ; deux capteurs de direction (écart toléré : 5\,\%) ; RTK licencié ; dashboard obligatoire & Le \og mode prudent \fg{} ne peut pas descendre sous 1,54\,m/s ; checklist automatique avant armement ; codes d'erreur affichés au pilote \\
Environnement technique & GNSS/RTK, liaison sol longue portée, calcul embarqué sur batterie & Repli sur retour au départ si la cible est perdue ; budget de latence partagé entre vision et planification \\
Phases de vie & Simulation, banc à terre, BlueBoat, lac, inspection, courses, transport, passation & Quatre modes logiciels : simulation, matériel dans la boucle, banc statique, déploiement \\
\bottomrule
\end{longtable}
\endgroup

\subsection{Missions}

\begin{tameoBox}[Mission principale]
Permettre au bateau IA de TAMEO de \textbf{percevoir les objets marins} qui
l'entourent et de \textbf{décider seul de les éviter ou de les atteindre},
afin d'accomplir les épreuves de la classe IA du MEBC \textbf{le plus vite
possible, sans contact et dans le respect du règlement}, sous la supervision
d'un pilote capable de reprendre la main à tout instant.
\end{tameoBox}

\begingroup
\footnotesize
\rowcolors{2}{tablerow}{white}
\begin{longtable}{L{1.75cm}L{2.75cm}L{3.0cm}L{3.4cm}L{2.9cm}}
\caption{Missions opérationnelles par épreuve}\label{tab:missions}\\
\toprule
\textbf{Épreuve} & \textbf{Connu / inconnu} & \textbf{Mission de détection} & \textbf{Mission d'évitement ou d'interception} & \textbf{Critère de succès} \\
\midrule
\endfirsthead
\toprule
\textbf{Épreuve} & \textbf{Connu / inconnu} & \textbf{Mission de détection} & \textbf{Mission d'évitement ou d'interception} & \textbf{Critère de succès} \\
\midrule
\endhead
Figure 8 & Deux bouées à position connue ; aucun obstacle & Confirmer les bouées, surveiller les abords & Suivre les points de passage GPS ; garder ses distances avec les bouées et les yachts & Temps minimal, sans contact \\
Slalom & Deux bouées fixes connues ; une bouée mobile entre elles & Détecter et suivre la bouée mobile & Éviter dynamiquement la bouée mobile en enchaînant les passages & Temps minimal, sans contact \\
Docking & Marque de virage à position inconnue ; départ et arrivée au quai & Trouver la marque, puis le quai et l'eau libre & Intercepter et contourner la marque, revenir, s'aligner et s'arrêter & À l'arrêt, étrave à 20\,cm du quai, sans toucher (sinon course annulée) \\
Sprint & Marque à moins de 0,4 NM, gisement inconnu ; autres bouées autour du stade & Rechercher et distinguer la bonne marque & Intercepter et contourner la marque sans sortir du stade & Temps minimal ; aucune sortie de zone \\
\bottomrule
\end{longtable}
\endgroup

\noindent Quatre \textbf{missions de soutien} complètent ces missions
opérationnelles :
\begin{itemize}
    \item \textbf{Qualifier} : démontrer au comité technique la checklist
    automatique, l'affichage des codes d'erreur, la reprise en main et la
    manœuvrabilité (inspection à terre, puis essai sur l'eau) ;
    \item \textbf{Sécuriser} : s'arrêter en fin de mission, actionneurs à
    0\textdegree{}, et exiger au moins deux actions physiques avant toute
    nouvelle mission ;
    \item \textbf{Rendre visible} : exposer en direct ou en différé ce que
    \og voit \fg{} et décide l'IA (dashboard) ;
    \item \textbf{Transmettre} : livrer un système documenté, testable en
    simulation et reprenable par les promotions suivantes.
\end{itemize}

% ===========================================================================
\section{Risques et prochaines étapes}
% ===========================================================================

L'analyse finale des risques du projet identifie, pour le SdI, des risques
encore ouverts dont la criticité initiale (probabilité $\times$ impact) est
élevée. Le tableau~\ref{tab:risques} les associe à une action prioritaire.

\begingroup
\footnotesize
\rowcolors{2}{tablerow}{white}
\begin{longtable}{L{4.6cm}L{1.0cm}L{9.45cm}}
\caption{Principaux risques ouverts et actions recommandées}\label{tab:risques}\\
\toprule
\textbf{Risque ouvert} & \textbf{Crit.} & \textbf{Action recommandée} \\
\midrule
\endfirsthead
\toprule
\textbf{Risque ouvert} & \textbf{Crit.} & \textbf{Action recommandée} \\
\midrule
\endhead
Bouée non détectée à longue distance (60\,m) & 12 & Mesurer la courbe portée/rappel sur l'eau, choisir entre le modèle précis et le modèle rapide \\
Latence de la vision trop élevée & 12 & Instrumenter la boucle complète, de l'image à la commande, et la comparer au budget de 300\,ms \\
Navigation autonome non validée sur le bateau réel & 12 & Recaler les limites dynamiques du DWA par essais sur l'eau \\
Faux positifs (reflets, vagues) & 9 & Fusionner les doublons et filtrer dans le temps avant de déclarer un obstacle \\
Écart entre images en ligne et scènes réelles & 9 & Enrichir le jeu de données avec des images tournées sur site \\
Collision si une bouée n'est pas détectée & 8 & Prédire la trajectoire de la bouée mobile (vitesse constante) et conserver les marges de sécurité \\
\bottomrule
\end{longtable}
\endgroup

\noindent\textbf{Feuille de route recommandée}, par ordre de priorité :
\begin{enumerate}
    \item \textbf{Mesurer avant d'optimiser} : portée, latence et taux de
    faux positifs sur l'eau, journalisés pour pouvoir rejouer les essais.
    \item \textbf{Anticiper la bouée mobile} : intégrer sa vitesse estimée dans
    le coût d'obstacle du DWA, puis valider l'anticipation en simulation.
    \item \textbf{Fiabiliser la liste d'obstacles} : fusion des doublons,
    filtrage temporel, recherche active quand la cible reste invisible.
    \item \textbf{Recaler la dynamique} du bateau de 5\,m et vérifier les
    conventions de signe de la vitesse de lacet.
    \item \textbf{Arbitrer la marge d'accostage} : l'arrêt actuel, à environ
    0,5\,m du quai, privilégie l'absence de contact ; le rapprocher des
    20\,cm réglementaires demande une distance au quai mesurée de façon fiable.
\end{enumerate}

\begin{tameoBox}[Message clé]
La chaîne de détection et d'interception existe de bout en bout, de la caméra
au safran, et elle est testable sans bateau grâce à la simulation. Sa
crédibilité dépend désormais de \textbf{mesures sur l'eau} plutôt que de
nouvelles fonctionnalités.
\end{tameoBox}

\vfill
{\scriptsize\color{tameoText}\noindent\textbf{Sources.} Rapport final et
dossier de réalisation du PIE (avril 2026) ; analyse des risques finale ;
manuel du module Vision ; règlement technique de la classe IA du MEBC
(version 2026.1) et description des épreuves ; cahier des charges du projet ;
code embarqué du bateau IA (dernière version consultée : juillet 2026). Les
résultats de la compétition ne figurent pas dans ces sources.\par}

\end{document}
